\section{Both interpretations and their composites}
\label{sec:interpretations}

The preceding argument supplies an actual numerical equality
characteristic and full intuitionistic logic. The rest of the arithmetic
translation uses the bounded-initializer construction, reproduced in
Appendix~\ref{app:bounded}. This section explains exactly what is inherited
and what must additionally be checked for the new primitive monus.

\subsection{Why graphs, rather than unrestricted output terms}
Tied-seed terms can still be bounded by a primitive-recursive unary
function of one input coordinate. That does not prevent a function from
being definable by a finite choice of terms. The bounded compiler supplies,
for every fixed primitive-recursive description $f$, source terms
\[
 \tau_f(b,\vec x),\qquad M_f(b),
\]
with proved bounds and coherence: whenever both $b$ and $c$ bound the
arguments, $\tau_f(b,\vec x)=\tau_f(c,\vec x)$. Neither $f$ nor a
description code is quantified over inside the source theory.

For a nonempty tuple $(x_1,\ldots,x_k)$, let $C_i(\vec x)$ say that $i$
is its first largest coordinate in the internally derived discrete
order. These are finitely many exhaustive, disjoint, quantifier-free
cases. The graph
\[
 \Phi_f(\vec x,z)=\bigvee_{i=1}^k
       \bigl(C_i(\vec x)\land z=\tau_f(x_i,\vec x)\bigr)
\]
is provably total and single-valued. Its outputs are explicit terms on
each case. Nullary functions have a single fixed value term.

Finite value cases are crucial for induction preservation. Expanding
every occurrence of a graph-defined function in a finite term gives a
finite family of quantifier-free guards and source value terms. Refining
these cases for every term in a quantifier-free formula gives an
equivalent source quantifier-free formula. This applies to implication
and negation as well, because the guards cover all inputs and the values
agree on overlaps. Source quantifier-free induction, already derived
from literal-equation induction and our equality characteristic,
therefore proves every translated iPRA induction instance.

An arbitrary existential graph would not establish this fact. The finite
term-value construction in the appendix is the relevant part of the
inherited proof.

\subsection{Arithmetic round trips}
Interpret the new $B+d$ language in iPRA by ordinary primitive recursion:
the initial value for $B_h$ is $S h(0,0,\vec p)$, its update is $h$, and
$d$ is ordinary truncated subtraction. The two $B$ schemas and the two
monus schemas hold. Interpret native equations as actual equality and
$Q,D$ by their ordinary readings. Coverage proves DI-case and equation
induction proves DI-progress. Every native rule is therefore sound.

Interpret iPRA back by the finite graphs above. The computation and
finite-case lemmas prove the defining equations and all induction
instances, and predecessor plus separation prove the successor boundary
laws. Both numerical interpretations have domain formula $x=x$ and use
actual equality.

On the iPRA side, external induction on each fixed primitive-recursive
description, with internal quantifier-free induction at a recursion
node, identifies its composite with the original function.

On the source side, first preserve the $d$-free constant-body recursor
$C$ and counter-body recursor $W$ by their equations and recurrence
uniqueness. This preserves $P=W(C(n,n))$. Next preserve primitive $d$
using its two equations, the preserved $P$, and finite quantifier-free
graph elimination. The constructed difference
$\widehat d$ of the inherited bridge has the same two recursive equations
as our primitive $d$. Literal-equation induction in its second argument
proves $\widehat d(x,y)=d(x,y)$; this source-level proof is included at the
end of the order appendix. Thus the new primitive is recovered, not just
some externally equivalent standard function. The equality characteristic
is then identified by its proved Boolean range and reflection. For each fixed staged
$B_h$, external induction on its earlier-stage body identifies the two
updates, and equation induction on the counter identifies their results.
When a graph-defined function occurs in a comparison, its finite
quantifier-free cases justify the induction in the original language.

\subsection{Formula round trips}
The gate and selector endpoint proofs identify each ordinary connective
with its native macro. The nontrivial native composite is
\[
 D_{x,y}A(x,y)\quad\dashv\vdash\quad
 \exists_*x\bigl(A(x,0)\Rightarrow_*\forall_*y\,A(x,y)\bigr).
\]
Forward, open the paired formula at a fresh witness. Under its zero
instance, QE provides every instance, giving the universal and then the
implication and existential macros. Reverse, eliminate the translated
existential, retain the matrix's zero instance, and use its implication
and universal at $Sz$. Case introduction reconstructs the original
paired formula. The unused-witness representative gives the analogous
composite for $Q$.

External structural induction on a fixed formula now proves both
formula composites, including arbitrary quantified formulas. All
quantifiers range over the original whole domain. The identity relation
$x=y$ defines each numerical isomorphism and preserves every operation,
relation, and translated formula. Ordinary falsity corresponds to the
native contradiction $0=1$, whose explosion was explicitly derived.

\begin{theorem}[The new native arithmetic basis]
The native Eq/$Q$/$D$ theory with B0, BS, M0, and MS, and exactly the
rules of Section~\ref{sec:rules}, is whole-domain bi-interpretable and
equiconsistent with full first-order iPRA. Both composite domain
isomorphisms are defined by actual equality.
\end{theorem}

The equiconsistency conclusion follows from effective proof
translations preserving contradiction. It is not a claim that either
theory proves its own consistency. Nor does it identify the implementation
cost of this basis with the smaller, independently measured $I$ checker.
